\section{What One Query Costs}
\label{sec:ch03-what-one-query-costs}

\index{statement count!single service}%
Delete the database file, start the service, and send one request. One request,
because the counter runs from the moment the process starts and never resets,
so the last number printed is what that request cost.

\begin{minted}{text}
POST /graphql HTTP/1.1
Host: localhost:5001
Content-Type: application/json

{"query":"{ sessions { title speaker { name } } }"}
\end{minted}

The answer is the one chapter~\ref{ch:hot-chocolate-zero} got. The console is
where the new information is.

\begin{minted}[fontsize=\footnotesize, breakanywhere]{text}
-- statement 1
SELECT "s"."Id", "s"."Abstract", "s"."DurationMinutes", "s"."SpeakerId", "s"."StartsAt", "s"."Title"
FROM "Sessions" AS "s"
ORDER BY "s"."StartsAt"
-- statement 2
SELECT "s"."Id", "s"."Bio", "s"."Name"
FROM "Speakers" AS "s"
WHERE "s"."Id" = @session_SpeakerId
LIMIT 1
-- statement 3
SELECT "s"."Id", "s"."Bio", "s"."Name"
FROM "Speakers" AS "s"
WHERE "s"."Id" = @session_SpeakerId
LIMIT 1
-- statement 4
SELECT "s"."Id", "s"."Bio", "s"."Name"
FROM "Speakers" AS "s"
WHERE "s"."Id" = @session_SpeakerId
LIMIT 1
-- statement 5
SELECT "s"."Id", "s"."Bio", "s"."Name"
FROM "Speakers" AS "s"
WHERE "s"."Id" = @session_SpeakerId
LIMIT 1
\end{minted}

\index{N+1@N+1!the shape of it}%
Five statements for four sessions. One to get the list, then one per session,
and statements two through five are identical to each other in every character
including the parameter name. That is the shape the chapter is named after:
one query for the collection, N for the rows in it, and the N is not a constant
you can look up. It is however many sessions the conference has. Four here.
Four hundred on a schedule anyone would pay to attend.

\index{Ada Fischer!gives two sessions}%
One of those four is pure waste. Ada Fischer gives the first two sessions on
the schedule, so statements two and three ask SQLite for the same row and get
the same answer back. The resolver has no way to know that. It is called with
one \code{Session} and returns one \code{Speaker}, and by the time it runs, the
list it belongs to is not something it can see.

\index{ordering!done in SQL}%
The good news is in statement one. \code{ORDER BY "s"."StartsAt"} means the
sort is happening in SQLite rather than in your process, which is the value
converter from section~\ref{sec:ch03-the-conference-in-sqlite} paying for
itself. Ask a database to sort and it sorts; ask it for everything and sort
afterwards and you have moved the whole table across a process boundary to
put four rows in order.

\index{statement count!with no nested field}%
One more measurement before the fix, because it isolates the cost. Restart, and
ask for the sessions without their speakers:

\begin{minted}{text}
{"query":"{ sessions { title } }"}
\end{minted}

That is one statement. Not five, not two: the speaker resolver is never called,
so it never asks. The extra four in the previous run were not the price of
loading the conference. They were the price of one field, and a client that
adds \code{speaker \{ name \}} to an existing query has quintupled the cost of
that query without touching your code or your schema.

In a controller you would have written that query yourself and seen what it
cost. In a graph the client writes it, you never see it, and the only thing you
control is what a single field costs when somebody asks for it.
